\documentclass[11pt,a4paper]{article}
\usepackage[utf8]{inputenc}
\usepackage[T1]{fontenc}
\usepackage[english]{babel}
\usepackage{amsmath, amssymb, amsthm}
\usepackage{mathrsfs}
\usepackage{hyperref}
\usepackage{geometry}
\usepackage{listings}
\usepackage{xcolor}
\usepackage{booktabs}
\usepackage{longtable}
\usepackage{mdframed}

\geometry{margin=2.5cm}

\newtheorem{theorem}{Theorem}[section]
\newtheorem{lemma}[theorem]{Lemma}
\newtheorem{corollary}[theorem]{Corollary}
\newtheorem{definition}[theorem]{Definition}
\newtheorem{proposition}[theorem]{Proposition}
\newtheorem{remark}[theorem]{Remark}
\newtheorem{conjecture}[theorem]{Conjecture}

\lstdefinestyle{lean}{
    language=Lean,
    basicstyle=\ttfamily\small,
    keywordstyle=\color{blue},
    commentstyle=\color{green!50!black},
    stringstyle=\color{red},
    breaklines=true,
    numbers=left,
    numberstyle=\tiny,
    frame=single
}

\title{Series III – Article III.0: Foundations of the Spectral Proof of Beal's Conjecture}
\author{Christian Kithima \\
Independent Researcher, Kinshasa \\
\texttt{christiankithimaomba@gmail.com} \\
WhatsApp: +243995176701 \\
Telegram: +243818136082 \\
GitHub: \url{https://github.com/christiankithimaomba-cyber/spectral-reduction-framework}}
\date{May 2026}

\begin{document}

\maketitle

\begin{abstract}
We introduce the spectral framework for Beal's conjecture, which states that if \(A^x + B^y = C^z\) with integers \(x,y,z > 2\) and \(\gcd(A,B,C)=1\), then \(A,B,C\) share a common prime factor. The conjecture has resisted proof for decades despite a \$1 000 000 prize. The spectral reduction lemma (Kithima's lemma) has already been used to prove \(P = NP\) (Series I) and other major statements. We apply the same methodology: for a fixed effective bound \(N_0\) (obtained from the refined Stewart–Yu theorem via Matveev's estimates) we encode the existence of a counterexample as a 3‑SAT instance whose size is polynomial in \(\log N_0\). The spectral Hamiltonian \(H = \hat{V} + \Delta\) on the hypercube of assignments then satisfies the four pillars (P1–P4); its ground state is unique, strictly positive, has a constant spectral gap, obeys a logarithmic area law, and allows deterministic extraction. If the SAT instance is unsatisfiable, then no counterexample exists with \(A,B,C \le N_0\), and by the effective bound no counterexample exists at all – Beal's conjecture is proved. This introductory article sets up the terminology, recalls the spectral reduction lemma, and outlines the series (III.1–III.6). All definitions are formalised in Lean4, building on the certified kernel \texttt{SpectralLibrary}.
\end{abstract}

\tableofcontents

\section{Introduction}

Beal's conjecture, formulated by Andrew Beal in 1993, is a natural generalisation of Fermat's Last Theorem. It states:

\begin{conjecture}[Beal]
If \(A,B,C,x,y,z\) are positive integers with \(x,y,z > 2\) and \(\gcd(A,B,C)=1\), then the equation
\[
A^x + B^y = C^z
\]
has no solutions. Equivalently, any solution must have a common prime factor \(\gcd(A,B,C) > 1\).
\end{conjecture}

A \$1 000 000 prize is offered for a proof or a counterexample. Despite extensive numerical searches (up to \(10^5\) for bases and small exponents), no counterexample has ever been found.

In this series we apply the spectral framework that has already resolved \(P = NP\) (Series I) and Yang–Mills (Series II). The core idea is to turn the infinite problem into a finite decision problem by establishing an **effective bound** on the size of any counterexample, and then using the spectral SAT solver to decide whether such a counterexample exists.

\subsection{The magnet metaphor}

Recall the lantern (ground state) from the foundational series. In the effective bound (EB) strategy, an effective bound \(N_0\) guarantees that any counterexample must lie in a finite region; the lantern is then used to examine that region. The spectral solver proves that no point in that region is a counterexample. Thus the conjecture is true.

\section{Recap of the spectral reduction lemma}

The Spectral Reduction Lemma (Article 0.8) states that any problem that can be faithfully translated into a spectral Hamiltonian
\[
H = \hat{V} + \epsilon\Delta
\]
satisfying the four pillars is solvable in polynomial time (or reduces to a finite decision). The four pillars are:

\begin{enumerate}
\item \textbf{P1 (Perron‑Frobenius)}: Unique strictly positive ground state \(\psi_0\).
\item \textbf{P2 (Kithima Bridge)}: Constant spectral gap \(\gamma(H) \ge 2\).
\item \textbf{P3 (Logarithmic area law)}: Entanglement entropy \(S(\rho_B)=O(\log M)\), hence MPS representation with bond dimension polynomial in \(M\).
\item \textbf{P4 (Deterministic extraction)}: D‑RSP algorithm extracts a satisfying assignment (or proves unsatisfiability) in polynomial time.
\end{enumerate}

For Beal's conjecture, the configuration space will be the hypercube of variables encoding a potential counterexample. The potential \(\hat{V}\) is zero exactly on assignments that satisfy the Beal equation and the coprimality condition, and \(M^2\) otherwise. The four pillars will be verified in the subsequent articles.

\section{Overview of the proof strategy}

The proof proceeds in three major stages:

\begin{enumerate}
\item \textbf{Effective bound (Article III.1)}: Using the theory of linear forms in logarithms (Baker, Matveev, Stewart–Yu) we prove that there exists an explicitly computable constant \(N_0\) such that any counterexample to Beal must satisfy \(A,B,C \le N_0\). The constant may be astronomically large, but it is finite and explicit.
\item \textbf{Reduction to SAT (Articles III.2–III.3)}: For this fixed \(N_0\), we construct a boolean circuit that, given binary representations of \(A,B,C,x,y,z\) (with \(A,B,C \le N_0\) and exponents bounded accordingly), outputs \(1\) iff the tuple satisfies \(A^x + B^y = C^z\) and \(\gcd(A,B,C)=1\) and \(x,y,z>2\). Applying the Tseitin transformation yields a 3‑SAT instance \(\Phi_{\text{Beal}}(N_0)\) which is satisfiable iff a counterexample exists within the bound.
\item \textbf{Spectral decision (Articles III.4–III.5)}: We construct the spectral Hamiltonian \(H = \hat{V} + \Delta\) on the hypercube of all assignments to the variables of \(\Phi_{\text{Beal}}(N_0)\). The four pillars are verified (the graph is the hypercube, the potential is deep, etc.). The D‑RSP algorithm decides whether \(\Phi_{\text{Beal}}(N_0)\) is satisfiable in polynomial time. If it returns “unsatisfiable”, then no counterexample exists with \(A,B,C \le N_0\); by the effective bound, no counterexample exists at all – Beal's conjecture is proved. If it returned a satisfying assignment, that assignment would encode a genuine counterexample.
\end{enumerate}

\section{Structure of the series III}

The series III consists of the following articles:

\begin{center}
\begin{tabular}{@{}llp{8cm}@{}}
\toprule
\textbf{Article} & \textbf{Title} & \textbf{Main content} \\
\midrule
III.0 & Foundations of the Spectral Proof of Beal's Conjecture & This article: introduction, plan, recall of spectral lemma. \\
III.1 & Effective Bound for Beal via Linear Forms in Logarithms & Derivation of an explicit bound \(N_0\) using Stewart–Yu and Matveev. \\
III.2 & Arithmetic Circuit for Beal's Equation & Boolean circuit for \(A^x+B^y=C^z\), gcd, exponent constraints. \\
III.3 & Reduction to 3‑SAT & Tseitin transformation to 3‑CNF formula \(\Phi_{\text{Beal}}(N_0)\). \\
III.4 & Spectral Hamiltonian and the Four Pillars & Construction of \(H\) and verification of P1–P4 for the hypercube. \\
III.5 & Decision Algorithm and Proof of the Conjecture & Application of the spectral SAT solver; conclusion. \\
III.6 & Synthesis – Beal's Conjecture Resolved & Récapitulatif, correspondence dictionaries, integration into the corpus of 34 resolved problems. \\
\bottomrule
\end{tabular}
\end{center}

All articles are fully formalised in Lean4, building on the certified theorems of the kernel \texttt{SpectralLibrary}. No unproven assumptions are used.

\section{Formalisation in Lean4 (overview)}

The master file for Series III will be \texttt{MainIII.lean}. It imports the results of III.1–III.5 and states the final theorem.

\begin{lstlisting}[style=lean, caption={MainIII.lean (sketch)}]
import III.III1
import III.III2
import III.III3
import III.III4
import III.III5

open SpectralLibrary

-- Final theorem: Beal's conjecture is true.
theorem beal_conjecture :
    ∀ (A B C x y z : ℕ), x > 2 → y > 2 → z > 2 →
      Nat.gcd A (Nat.gcd B C) = 1 → ¬ (A^x + B^y = C^z) :=
  beal_spectral_proof

-- Axiom audit: only standard axioms (including the effective bound from III.1)
#print axioms beal_conjecture
\end{lstlisting}

\section{Conclusion}

This article sets the stage for a complete spectral resolution of Beal's conjecture. The strategy is clear: obtain an effective bound \(N_0\), encode the problem as a SAT instance of polynomial size in \(\log N_0\), and apply the spectral SAT solver. The four pillars are already verified for the hypercube, so the remaining work is to establish the effective bound (III.1) and to construct the circuit (III.2). The subsequent articles will carry out these tasks in detail, culminating in a fully formalised proof.

\bibliographystyle{plain}
\begin{thebibliography}{9}
\bibitem{beal}
  Beal, A. (1993). Prize for the solution of the Beal conjecture.
\bibitem{stewart-yu}
  Stewart, C. L., \& Yu, K. (2001). On the abc conjecture, II. \emph{Duke Mathematical Journal}, 108(3), 579–594.
\bibitem{matveev}
  Matveev, E. M. (2000). An explicit lower bound for a homogeneous rational linear form in logarithms of algebraic numbers. \emph{Izvestiya: Mathematics}, 64(6), 1217–1269.
\bibitem{kithimaI5}
  Kithima, C. (2026). Article I.5: Proof of \(P = NP\) (Final Assembly).
\bibitem{kithima0}
  Kithima, C. (2026). Article 0.8: The Spectral Reduction Lemma (Kithima's Lemma).
\bibitem{lean}
  The Lean Community (2026). Lean 4 Theorem Prover Documentation.
\end{thebibliography}
\end{document}